% SUBMISSION: {"kind": "research_attempt", "category": "economics", "problem_id": "L13-76", "model": {"provider": "OpenAI", "version": "Codex, GPT-6 family", "version_uncertainty": "The serving context exposed the GPT-6 family. Exact serving revisions were not exposed, including for other materially involved Codex agents; no revision is inferred from a folder.", "reasoning": "Not exposed by the session.", "generated_on": "2026-10-11"}, "other_models": "Separate Codex agents materially constructed or reviewed the finite-market and entry primitives, public-state closure, mixing/bias and replacement proofs, robustness, examples, original scope, sources and presentation. Their exact serving revisions and reasoning configurations were not exposed. No other material model family is confirmed.", "tools": "File and Git/GitHub tools; primary-source web and PDF inspection; LaTeX compilation/rendered-page inspection as recorded in the production report. Python/Node repository and site checks are publication checks. No Lean or other theorem prover, computer-algebra proof, simulation or executed pricing/linear-program/quantifier-elimination solver is claimed.", "target": "L13-76, Algorithmic pricing conduct: identify conditions for an economically material persistent quantity-weighted expected price gap over the same-environment selected one-period optimal benchmark, surviving every allowed unilateral algorithm replacement; include demand, exploration and entrant perturbations, distinguish reciprocal punishment from unilateral commitment, and state architecture/stationarity dependencies.", "scope": "Complete proof claim for the original declared-algorithm-class target in a fixed finite price grid inside the common compact interval, a reconstructible finite public market/controller state, finite entry and shock primitives, and stationary finite public Markov baseline controllers. All joint physical transition kernels satisfy a common full-support Doeblin minorization and total quantity has a positive floor. Replacements range over every measurable randomized history program on the grid using the allowed public and own-profit information; they are not restricted to stationary controllers. The same physical demand/cost/shock, entrant response and entry-opportunity primitives and initial law generate the selected verified one-period Nash benchmark at its actual prices. The manuscript characterizes price/profit limits, persistence and replacement survival by finite invariant/Poisson-bias/slack conditions, gives quantitative and exact semialgebraic perturbation criteria, and separates reciprocal continuation reactions, entry effects, unilateral commitment and exploration fragility. No universal statement for arbitrary learners, hidden predictive memory, all continuous prices or unrestricted nonstationary markets is claimed.", "human_contributions": "JiHo Bae supplied the requested range, nonduplication/preservation, readable and accurately cited proof, and GitHub/PR instructions. No significant human mathematical contribution, substantive proof correction or independent human proof checking is confirmed. OpenAI Codex is the primary mathematical generating attribution; JiHo Bae is credited as submitter and prompting contributor.", "verification": "Separate Codex agents reviewed mathematical, original-scope, source and symbolic-example components. These are model-assisted reviews. Accompanying final records identify the complete frozen source and actual audit timing/outcomes; this declaration does not assert an unperformed or successful audit. No independent human proof verification, formalization, computer-algebra proof, simulation or executed solver is claimed. Compiler/rendering and repository/site checks concern production/software, not mathematical correctness or novelty.", "reproducibility": "Literal Prompt A was applied on 2026-10-11 at 05:39:24 KST before prospective full-proof construction; prior activity was candidate/model screening. The full registered original, actual policy and filled prompt/input hashes were retained. A later actual comparison confirmed those bytes identical at official baseline b1f633121074e16579fa7fd05314e5f504d29bae. The final reproducibility record identifies the frozen full proposed source, actual filled Prompt B inputs, timing and outcome before submission. Public source, bibliography, review receipts and hashes accompany the release. Exact serving revisions/reasoning and full original generation transcript are unavailable; no private transcript or secret is published."}
% FULL_SOLUTION_SCOPE: {"target": "L13-76, Algorithmic pricing conduct: identify conditions for an economically material persistent quantity-weighted expected price gap over the same-environment selected one-period optimal benchmark, surviving every allowed unilateral algorithm replacement; include demand, exploration and entrant perturbations, distinguish reciprocal punishment from unilateral commitment, and state architecture/stationarity dependencies.", "result": "Lemma 1 proves mixing, invariant averages and the normalized Poisson bias. Theorem 2 proves the exact necessary-and-sufficient finite slack test for all allowed history-dependent randomized replacements, the span(h)/T finite-horizon bound and profitable one-state witness. The stationary price formula and gap condition in Section 4 give exact persistence against the same-environment benchmark. Proposition 3 gives quantitative perturbation bounds; Corollary 4 gives exact semialgebraic global-robustness recognition and algebraic witnesses. Proposition 5 supplies the entry/reset threshold and demand/entry formulas; Proposition 6 proves exploration-fragility with an explicit profitable replacement. Section 7 distinguishes punishment, commitment and architecture/stationarity conditions.", "coverage": "The original permits a declared algorithm class and finite/compact program menus. Within the stated finite-grid/public-state stationary architecture, the proof covers all allowed finite firms, states, grids, finite shock/entry outcomes, initial laws, nonzero quantity bounds and reset parameters including delta=1; all grid prices and arbitrary infinite-history randomized replacements with the allowed observations; both baseline and same-environment selected static benchmark under their own induced state laws; equality at the prespecified material threshold; every finite nonnegative slack or negative-slack witness; declared demand, exploration and entrant perturbations preserving the architecture/quantity/minorization conditions; exact survival versus approximate gain and support/exploration boundaries; finite or stationary-augmented architecture and the limits of nonstationary application; and the original distinction between observed high prices, replacement equilibrium, reciprocal punishment and unilateral commitment. Entry costs/entrant behavior are retained rather than silently removed. Used baseline actions have zero slack, so no fictitious strictly negative best-deviation margin is used to claim general open robustness.", "dependencies": "Carl Shapiro and Ali Yurukoglu, Trends in Competition in the United States: What Does the Evidence Show?, author draft dated 24 August 2024, Section 6, printed p.37, fourth research direction: motivating pricing/entry/conduct agenda, not the finite-controller theorem. Saugata Basu, Richard Pollack and Marie-Francoise Roy, Algorithms in Real Algebraic Geometry, second edition (2006), Sections 13.1 and 14.2: exact real-algebraic sampling and quantifier elimination for the explicitly formulated finite polynomial systems over effectively encoded rational/algebraic data. No polynomial complexity for general quantifier elimination is claimed. Mixing, invariant/bias identities, history-policy telescoping and one-state replacement necessity, perturbation bounds, benchmark best-response checks and examples are supplied by proof; no imported generic average-reward-MDP or unverified Nash-existence theorem is required.", "human_verification": "None confirmed. Separate Codex-agent checks are model-assisted reviews, not independent human proof checking.", "unverified": "The entire proof has not been independently checked by a human. No Lean/formalization, computer-algebra proof, simulation or execution of a pricing/LP/QE solver is claimed. Symbolic boundary/example checks and cited-hypothesis reviews are model assisted; accompanying final records identify their actual frozen-source outcomes. No missing mathematical obligation is asserted within the expressly declared original-target class. Solved/100 percent expresses a whole-scope proof claim, not confidence, mathematical certification or independent validation.", "novelty": "The earlier L13-76 partial attempt with no-entry grim punishment and absorbing-experimentation limitations is preserved and distinguished; it is not repackaged as the new full result. The proposed core is the exact reduction of the original unrestricted allowed-history replacement and persistent-price conditions to finite economic primitives under uniform recurrence, with same-environment entry benchmark and complete perturbation/mechanism analysis. Primary-source/context and previous-attempt checks were performed. No exhaustive literature/equivalence search or established worldwide novelty claim is made; duplicate-clearance evidence is recorded separately."}
% ECONOMICS_PROBLEM: {"id":"L13-76","title":"Algorithmic pricing conduct","jel_code":"L13","source_id":"OP-0398"}
% FIRST_POSTED: 2026-10-11
% ATTEMPT_STATUS: solved
% ENTRY_KIND: research_attempt
\documentclass[11pt,a4paper]{article}
\usepackage[T1]{fontenc}
\usepackage[utf8]{inputenc}
\usepackage{lmodern,amsmath,amssymb,amsthm,mathtools,microtype}
\usepackage[margin=24mm]{geometry}
\usepackage{xurl,hyperref,enumitem}
\hypersetup{colorlinks=true,linkcolor=blue,citecolor=blue,urlcolor=blue,
 pdftitle={Algorithmic pricing with entry: exact persistence and replacement tests}}
\setlength{\parindent}{1.1em}
\setlength{\parskip}{3pt}
\setlength{\emergencystretch}{2em}
\setlist{nosep,leftmargin=1.8em}
\newtheorem{theorem}{Theorem}
\newtheorem{lemma}[theorem]{Lemma}
\newtheorem{proposition}[theorem]{Proposition}
\newtheorem{corollary}[theorem]{Corollary}
\newcommand{\R}{\mathbb R}
\newcommand{\E}{\mathbb E}
\newcommand{\Prob}{\mathbb P}
\newcommand{\ind}{\mathbf 1}
\newcommand{\spn}{\operatorname{sp}}
\newcommand{\TV}{\operatorname{TV}}
\newcommand{\supp}{\operatorname{supp}}
\title{Algorithmic pricing with entry:\\exact persistence and replacement tests}
\author{OpenAI Codex (GPT-6 family)\\
\small Exact serving revision and reasoning setting not exposed\\
\small Submitted by JiHo Bae}
\date{11 October 2026}
\begin{document}
\maketitle

\paragraph{Submission provenance.}
The complete mathematical argument, scope and contribution disclosures are
retained from the \href{https://github.com/jbaelaw/pricing-program-stability-proof/blob/0358e162cfacff2b7c001955a9ca0076244be99c/paper/Proof.tex}{immutable public source} and
\href{https://github.com/jbaelaw/pricing-program-stability-proof/blob/0358e162cfacff2b7c001955a9ca0076244be99c/paper/Proof.pdf}{proof PDF}. This catalogue wrapper records publication
provenance. Community review remains pending.

\begin{abstract}
For a declared class of finite public-state pricing programs, we characterize
when an above-benchmark quantity-weighted price persists and survives replacement
of any one firm's entire algorithm. The replacement may use arbitrary public
history, past own profits and private randomization. A common full-support reset
reduces this infinite-history condition to finitely many Poisson inequalities;
a violated inequality produces a profitable one-state replacement. The benchmark
uses the same demand, entry opportunities and entrant responses. Explicit
perturbation bounds distinguish robust price gaps from fragile replacement
incentives. A market with endogenous finite-cost entry gives a positive stable
gap and an exact reset/entry threshold, while arbitrarily small exploration can
destroy exact stability without eliminating the gap.
\end{abstract}

\section{The target and the market architecture}
Problem L13-76 \cite{catalogue} asks for conditions under which a quantity-weighted
price gap survives increasing horizons and unilateral algorithm replacement.
It also asks for demand, exploration and entry perturbations, a distinction
between reciprocal punishment and unilateral commitment, and the dependence on
architecture and stationarity. Shapiro and Yurukoglu
\cite[Section 6, p.37, fourth research direction]{shapiro} motivate this pricing
and entry agenda. They do not supply the finite-controller theorem below.

We answer all these questions for the following fixed class, using the original
permission to declare a program architecture. There are $N\geq2$ incumbents,
a finite nonempty grid $B\subset[0,p_{\max}]$, and a finite nonempty public
state space $S$. Every comparison starts from the same fixed law $\alpha$.
One infinite market and one infinite policy are evaluated at horizons
$T=1,2,\ldots$; neither is changed with $T$.

The state records physical demand and observed cost variables, any persistent
entrant status, and the finite registers of specified controllers. Registers
and their updates are known and reconstructible from initial public conditions,
realized prices and quantities, and declared public physical shocks. If a
firm's original register is needed after its program is replaced, its shadow
register continues the same known update from the actual public observations.
There are no additional program messages or hidden predictive registers.

For each $(s,\boldsymbol b)\in S\times B^N$, a fixed finite-outcome law
specifies entry-cost draws, entry and exit decisions, entrant prices, demand
shocks, quantities, observed costs and the next public state. Entrant programs
and the timing of entry relative to incumbent prices are primitives of this
law. In particular, a fixed entrant program may respond to the actual current
incumbent prices. Every baseline, replacement and benchmark uses this same
response law and the same shock and entry opportunities. Conditional on current
state and current incumbent prices, the law has no further history dependence.
Fresh private entry draws are permitted if they leave no omitted predictive
information. Own profit observations are functions of public prices, quantities,
observed costs and state, or give no additional information about future
primitives conditional on them.

For each realized outcome, let quantities be nonnegative and define total
quantity and expenditure, including entrants, by
\[
 Q=\sum_jq_j,\qquad X=\sum_jp_jq_j,
 \qquad 0<q_{\min}\leq Q\leq q_{\max},\qquad 0\leq X\leq p_{\max}Q.
\]
Incumbent profit is the market payoff $p_iq_i-C_i(q_i)$, bounded on the finite
outcome tables. These prices and profits are derived from demand, allocation,
cost and entry primitives, rather than unrelated reward labels.
Write $L(s'\mid s,\boldsymbol b)$ for the resulting next-state law.
Fix $0<\delta\leq1$ and a full-support probability vector $\nu$ and require
\begin{equation}\label{eq:minorization}
 L(s'\mid s,\boldsymbol b)\geq\delta\nu(s')
 \quad\text{for every }s,s'\in S,\ \boldsymbol b\in B^N.
\end{equation}
An independent public reset with probability $\delta$ implements this assumption.
It applies to every joint price vector, including benchmark and deviation prices.

The fixed baseline program of firm $i$ draws independently from
$a_i(s)\in\Delta(B)$ each period. Registers can encode finite learning and
exploration rules; stationary randomization may also be present. A permitted
replacement is any measurable program using the entire allowed public history,
own past profits and its independent randomization to select prices in $B$.
Replacement is not restricted to finite memory or stationary behavior.
The finite grid is the declared pricing architecture inside the original common
compact interval. We make no claim about replacements allowed every real price
in that interval.

\paragraph{The benchmark.}
At each state, a fixed selected mixed profile $a^B(s)$ is supplied for the
one-period pricing game on the same grid with the same entrant law, costs and
shocks. Its stage best-response inequalities must be verified as part of the
input: every pure price has profit no greater than the selected mixed payoff
against the other selected prices. No existence or selection follows merely
from declaring a compact set. The explicit market in Section~\ref{sec:example}
verifies its selection directly. Benchmark state and register updates use
the benchmark's actual observed outcomes, so its stationary distribution need
not equal the baseline distribution.

\section{Finite-state closure and average rewards}
Against baseline rivals, integrate their independent prices and the fixed
entrant and demand law to define
\[
 r_i(s,b)=\E[p_iq_i-C_i(q_i)\mid s,b],\qquad
 P_i(s,s'\mid b)=\Prob(S_{t+1}=s'\mid S_t=s,b_i=b).
\]
The declared observability and outcome law imply that these are the correct
conditional quantities after every allowed replacement history. Rival registers
update from the replacement's actual public prices, not from a counterfactual
baseline path. Each $P_i(\cdot\mid s,b)$ satisfies \eqref{eq:minorization}.
The baseline chain and reward satisfy
\[
 P(s,s')=\sum_ba_i(s,b)P_i(s,s'\mid b),\qquad
 r_i(s)=\sum_ba_i(s,b)r_i(s,b).
\]
The same $P$ results for every $i$. Let $P^B$ be the chain obtained by integrating
$L$ under $a^B$. It also satisfies \eqref{eq:minorization}.

For a vector $v$ write $\spn(v)=\max_sv(s)-\min_sv(s)$, and for probability
vectors use total variation $\|\lambda-\lambda'\|_{\TV}
=\frac12\sum_s|\lambda(s)-\lambda'(s)|$.

\begin{lemma}[Mixing, averages and the Poisson equation]\label{lem:poisson}
Every stochastic matrix satisfying \eqref{eq:minorization} has a unique
stationary law $\mu$, with $\mu(s)\geq\delta\nu(s)>0$ and
\[
 \|\lambda P^t-\mu\|_{\TV}\leq(1-\delta)^t,
 \qquad \spn(Pv)\leq(1-\delta)\spn(v).
\]
For a reward vector $r$, put $g=\mu r$. Then
\begin{equation}\label{eq:averagebound}
 \left|T^{-1}\sum_{t=0}^{T-1}\alpha P^tr-g\right|
 \leq\frac{\spn(r)}{\delta T}.
\end{equation}
For any fixed $s_0$, there is a unique vector $h$ with $h(s_0)=0$ satisfying
\begin{equation}\label{eq:poisson}
 g\boldsymbol1+h=r+Ph,
 \qquad \spn(h)\leq\frac{\spn(r)}\delta.
\end{equation}
\end{lemma}
\begin{proof}
If $\delta<1$, write
$P=\delta\boldsymbol1\nu^\top+(1-\delta)R$ with $R$ stochastic.
The common first term cancels in differences of probability vectors and is
constant in differences of function values. Stochastic averaging cannot
increase total variation or span, giving both contractions. For $\delta=1$
every row equals $\nu^\top$, and the conclusions follow directly.

Consecutive iterates of any probability vector are a geometric Cauchy sequence.
Its limit is stationary; contraction gives uniqueness and the displayed mixing
bound. Stationarity and minorization give $\mu(s)\geq\delta\nu(s)$.
The expectation inequality
$|\lambda v-\lambda'v|\leq\spn(v)\|\lambda-\lambda'\|_{\TV}$ then gives
\eqref{eq:averagebound} by summing the geometric series.

The series $\widetilde h=\sum_{t\geq0}(P^tr-g\boldsymbol1)$ converges
uniformly by the same bound. Subtracting its image under $P$ telescopes to
$(I-P)\widetilde h=r-g\boldsymbol1$. Normalize by subtracting
$\widetilde h(s_0)$. Span contraction bounds its span by
$\sum_{t\geq0}(1-\delta)^t\spn(r)=\spn(r)/\delta$.
Two normalized solutions differ by a harmonic vector. Its span is at most
$(1-\delta)$ times itself, so it is constant and normalization makes it zero.
The argument includes singleton state spaces and $\delta=1$.
\end{proof}

\section{An exact test against every history-based replacement}
For each firm, compute $g_i=\mu r_i$ and its normalized bias $h_i$ from
\eqref{eq:poisson}. Define the finite deviation slacks
\begin{equation}\label{eq:slack}
 d_i(s,b)=g_i+h_i(s)-r_i(s,b)-P_i(s,\cdot\mid b)h_i.
\end{equation}
The baseline identity gives
\begin{equation}\label{eq:supportzero}
 \sum_ba_i(s,b)d_i(s,b)=0\qquad\text{for every }i,s.
\end{equation}

\begin{theorem}[Replacement survival]\label{thm:replacement}
The original average-profit condition
\[
 \limsup_{T\to\infty}
 \{U_{i,T}(a'_i,a_{-i})-U_{i,T}(a)\}\leq0
\]
holds for every firm and every allowed infinite-history randomized replacement
if and only if $d_i(s,b)\geq0$ for every $i,s,b$.
When it holds, every such replacement satisfies the uniform finite-horizon bound
\begin{equation}\label{eq:deviationbound}
 U_{i,T}(a'_i,a_{-i})-U_{i,T}(a)
 \leq\frac{\spn(h_i)}T\leq\frac{\spn(r_i)}{\delta T}.
\end{equation}
If $d_i(s_*,b_*)<0$, the stationary program that chooses $b_*$ only at $s_*$
and uses the baseline elsewhere is a profitable replacement.
\end{theorem}
\begin{proof}
Suppose all slacks are nonnegative. Conditional on the replacement's full
allowed history and current price, \eqref{eq:slack} gives
$\E[r_{it}+h_i(S_{t+1})\mid\text{current history}]
\leq g_i+h_i(S_t)$, also after averaging current private randomization.
Summing and taking expectations gives
\[
 U_{i,T}(a'_i,a_{-i})\leq g_i+
 \frac{\E_\alpha h_i(S_0)-\E^{a'_i}h_i(S_T)}T.
\]
The baseline Poisson equation gives equality with its own terminal law.
The initial terms cancel because the initial law is identical. The difference
of terminal expectations is at most $\spn(h_i)$, proving
\eqref{eq:deviationbound} and the required limsup. This bounds arbitrary
history dependence, rather than assuming a Markov replacement.

Conversely, consider the stated one-state switch. Its chain $P'$ inherits
\eqref{eq:minorization}, so its stationary law satisfies
$\mu'(s_*)\geq\delta\nu(s_*)>0$. By \eqref{eq:supportzero}, its expected
slack is zero away from $s_*$ and equals $d_i(s_*,b_*)$ there. Average the
slack identity under $\mu'$. The bias terms cancel by stationarity, yielding
\[
 g_i'-g_i=-\mu'(s_*)d_i(s_*,b_*)
 \geq\delta\nu(s_*)[-d_i(s_*,b_*)]>0.
\]
Lemma~\ref{lem:poisson} makes both finite-horizon average profits converge
to their gains from the specified initial law. The replacement therefore
violates the required limsup. This proves necessity and gives a concrete
finite public-state witness.
\end{proof}

Thus survival is a finite linear-algebra test, not a tautological optimization
over algorithms. With rational primitive tables and a rational selected
benchmark, stationary laws, normalized biases and slacks use exact rational
linear equations and comparisons. Exact decision claims also require the
threshold and all comparison constants to have rational or real-algebraic
encodings; the theoretical criteria allow arbitrary real constants.
By \eqref{eq:supportzero}, survival also
forces every price in the baseline support to have zero slack. A strictly
positive margin on all prices is impossible.

\section{The quantity-weighted persistence condition}
Let $x(s),q(s)$ be baseline expected expenditure and total quantity at state
$s$, and let $x^B,q^B$ be the corresponding benchmark vectors derived from
the same primitive outcome law. Let $\mu^B$ be the stationary law of $P^B$.
The exact finite-horizon ratios in the original problem are
\[
 \bar p_T=\frac{\sum_{t<T}\alpha P^tx}{\sum_{t<T}\alpha P^tq},
 \qquad p_T^B=\frac{\sum_{t<T}\alpha(P^B)^tx^B}
                       {\sum_{t<T}\alpha(P^B)^tq^B}.
\]
Every denominator is at least $q_{\min}T$. Lemma~\ref{lem:poisson} gives
\begin{equation}\label{eq:pricelimit}
 \bar p_T\longrightarrow\frac{\mu x}{\mu q},\qquad
 p_T^B\longrightarrow\frac{\mu^Bx^B}{\mu^Bq^B}.
\end{equation}
Consequently the prespecified material gap $\varepsilon>0$ persists exactly when
\begin{equation}\label{eq:gap}
 (\mu x)(\mu^Bq^B)-(\mu^Bx^B)(\mu q)
 \geq\varepsilon(\mu q)(\mu^Bq^B).
\end{equation}
Equality is allowed by the original liminf definition. These are ratios of
expected cumulative amounts, not expected ratios of random prices.

For completeness, put $A_T=T^{-1}\sum_{t<T}\alpha P^tx$,
$B_T=T^{-1}\sum_{t<T}\alpha P^tq$, $A=\mu x$ and $B_0=\mu q$.
Then $B_T,B_0\geq q_{\min}$, $A/B_0\leq p_{\max}$ and
\[
 \left|\frac{A_T}{B_T}-\frac A{B_0}\right|
 \leq\frac{|A_T-A|+p_{\max}|B_T-B_0|}{q_{\min}}
 \leq\frac{\spn(x)+p_{\max}\spn(q)}{\delta q_{\min}T}.
\]
The same bound for the benchmark quantifies convergence of the measured gap.
Together, \eqref{eq:gap} and the nonnegative slacks are necessary and sufficient
for the original bundled persistence and replacement requirements in this class.

\section{Demand, exploration and entry perturbations}
Perturbations change the declared finite demand and cost tables, entry-cost
probabilities, entrant responses or baseline exploration probabilities. Recompute
the induced baseline and unilateral tables; exploration need not leave their
transition matrices fixed. The following bounds compare any two such environments
on the same states and grids with a common minorization $\delta\nu$ and quantity
floor $q_{\min}$. Primes denote perturbed objects.

For one firm put
\[
 \eta=\max_s\|P'(s,\cdot)-P(s,\cdot)\|_{\TV},\quad
 e_r=\|r'-r\|_\infty,\quad H=\spn(r)/\delta,
\]
and let $e_d$ be the maximum change of a unilateral reward and $\eta_d$ the
maximum total-variation change of a unilateral transition row, over all states
and prices.

\begin{proposition}[Explicit continuity bounds]\label{prop:robust}
One has
\begin{align*}
 \|\mu'-\mu\|_{\TV}&\leq\eta/\delta,\\
 |g'-g|&\leq e_r+\spn(r)\eta/\delta,\\
 \spn(h'-h)&\leq[\spn(r'-r)+2\eta H]/\delta,\\
 \max_{s,b}|d_i'(s,b)-d_i(s,b)|
 &\leq e_r+\spn(r)\eta/\delta
       +[\spn(r'-r)+2\eta H]/\delta+e_d+\eta_dH.
\end{align*}
For $R=\mu x/(\mu q)$, if $\|x'-x\|_\infty\leq e_x$ and
$\|q'-q\|_\infty\leq e_q$, then
\begin{equation}\label{eq:priceperturb}
 |R'-R|\leq C_A:=\frac{e_x+\spn(x)\eta/\delta
       +p_{\max}[e_q+\spn(q)\eta/\delta]}{q_{\min}}.
\end{equation}
The analogous benchmark bound $C_B$ gives a gap change at most $C_A+C_B$.
\end{proposition}
\begin{proof}
Stationarity and contraction give
$\|\mu'-\mu\|_{\TV}\leq\eta+(1-\delta)\|\mu'-\mu\|_{\TV}$.
The gain bound follows from
$g'-g=\mu'(r'-r)+(\mu'-\mu)r$ and the expectation/span inequality.
Set $u=h'-h$ and subtract the Poisson equations:
\[
 u=(r'-r)-(g'-g)\boldsymbol1+P'u+(P'-P)h.
\]
The constant disappears under span, while each coordinate of $(P'-P)h$ has
absolute value at most $\eta\spn(h)\leq\eta H$. Span contraction therefore
gives $\delta\spn(u)\leq\spn(r'-r)+2\eta H$.
Finally subtract the slack identities:
\[
 d_i'-d_i=(g'-g)+u(s)-P_i'(s,\cdot\mid b)u
 -(r_i'(s,b)-r_i(s,b))-(P_i'-P_i)(s,\cdot\mid b)h.
\]
A coordinate minus a probability average is bounded by $\spn(u)$; the last
term is bounded by $\eta_dH$. This proves the slack bound, independently of
the bias normalization.

The numerator averages differ by at most $e_x+\spn(x)\eta/\delta$, and the
denominator averages by at most $e_q+\spn(q)\eta/\delta$. Use
\[
 R'-R=\frac{(\mu'x'-\mu x)-R(\mu'q'-\mu q)}{\mu'q'},
 \qquad 0\leq R\leq p_{\max},\quad\mu'q'\geq q_{\min},
\]
to obtain \eqref{eq:priceperturb}. The benchmark is covered because its
actual induced chain satisfies the same assumptions.
\end{proof}

A strict original gap $R-R^B\geq\varepsilon+\Gamma$, $\Gamma>0$, survives
every admissible perturbation with $C_A+C_B\leq\Gamma$. Changes in a selected
stage benchmark must be included in $C_B$; an arbitrary discontinuous equilibrium
selection does not inherit a continuity claim. A gap exactly at the threshold
has no automatic neighborhood guarantee.

Exact replacement survival requires a different argument. For a deterministic
baseline, its chosen action has zero slack identically. If every other action
has slack at least $\gamma>0$, perturbations preserving those chosen actions
and with the slack bound at most $\gamma$ preserve survival. Structurally zero
slacks, for example an inactive firm's irrelevant prices, may also be retained
if their zero identities are preserved; otherwise they must be recertified.
For a mixed baseline, all supported slacks must remain zero as well as all
other slacks nonnegative. Continuity alone does not preserve these equalities.
A full-support baseline survives only when every price has zero slack.

There is also an exact exploration failure. Suppose $d_i(s,b)>0$ at a surviving
profile. Keep rivals and market primitives fixed, and at state $s$ change only
firm $i$'s distribution to $(1-t)a_i(s)+t\delta_b$, $0<t\leq1$.
The old bias identity, averaged under the new stationary law $\mu^t$, gives
\[
 g_i^t-g_i=-t\mu^t(s)d_i(s,b)<0.
\]
Reverting to the original program is an allowed profitable replacement.
Thus adding a strictly inferior exploration action destroys survival for every
positive rate, although a strict price gap persists for sufficiently small
rates by Proposition~\ref{prop:robust}. Exploration among zero-slack actions
requires a separate test; no universal harm claim is made.

\subsection*{Exact robustness over an encoded parameter family}
Let $\Xi$ be a nonempty compact semialgebraic uncertainty set specified by
rational polynomial relations, with a rationally or real-algebraically encoded
material threshold $\varepsilon$ and bounds. Suppose its finite outcome probabilities,
quantities, costs, entry rules, controller probabilities and induced tables
are polynomial or rational functions of $\xi\in\Xi$, with every rational
denominator nonzero and of declared sign. The common minorization and quantity
floor are part of the family. Require a specified semialgebraic graph with
one selected benchmark profile for each parameter. Its stage best-response
inequalities are also verified; inequalities alone do not choose among equilibria.

\begin{corollary}[Finite recognition of global robustness]\label{cor:qe}
For this encoded family, whether persistence and survival hold at every
$\xi\in\Xi$ is decidable by quantifier elimination over the reals. If they
fail, a real-algebraic parameter witness and a failed price condition or
profitable one-state replacement can be extracted. No polynomial running-time
claim is made.
\end{corollary}
\begin{proof}
Introduce variables for $\mu,\mu^B,g_i,h_i$ and the uniquely selected benchmark
and any auxiliaries in its graph. Impose stationarity, probability normalization
and nonnegativity, the normalized Poisson equations, the benchmark graph and
its stage best-response inequalities. Add the finitely many
$d_i(s,b)\geq0$ and \eqref{eq:gap}. Clearing rational denominators according
to their known signs makes this a finite polynomial system. The gain and
normalized bias and both invariant laws are unique by
Lemma~\ref{lem:poisson}. Thus feasibility is equivalent to the exact
infinite-history property, not a relaxation of it.

Global robustness is the real first-order sentence asserting that for every
$\xi\in\Xi$ these certificate variables exist. The quantifier-elimination
and algebraic sampling algorithms in \cite[Sections 13.1 and 14.2]{bpr}
decide it and produce a sample from a nonempty semialgebraic failure set.
At that parameter, a negative slack yields Theorem~\ref{thm:replacement}'s
one-state program, or \eqref{eq:gap} fails. If the supplied benchmark itself
fails a best-response requirement, that is an invalid-input witness instead
of a pricing counterexample. The statement requires exact rational or
real-algebraic input encoding; it does not apply to unspecified real functions
or an arbitrary equilibrium correspondence.
\end{proof}

\section{A market with entry and an exact threshold}\label{sec:example}
Two incumbents have marginal cost zero and permitted prices $B=\{0,3/4,1\}$.
A unit mass of consumers values a unit at one, buys at the lowest price including
at one, and splits equally among tied lowest-price sellers. Incumbents observe
the public state $C$ or $P$ and simultaneously post prices. A new one-period
entrant then observes those prices and its independently drawn fixed entry cost
$F$, with $\Prob(F=0)=z$ and $\Prob(F=2)=1-z$, $0\leq z\leq1$.
Its marginal cost is $3/4$ and it uses the same price grid. The selected entrant
program enters exactly when $F=0$ and then quotes one. It exits after demand
is served.

This program is optimal at every incumbent price pair. Against $(1,1)$,
price one gives operating profit $(1-3/4)/3=1/12$, whereas price $3/4$
has zero margin and price zero has negative margin. If an incumbent prices
at most $3/4$, no grid price can earn positive operating profit; entry at
zero fixed cost and quote one is weakly optimal and is the selected tie rule.
Thus zero-cost entrants may enter but make no sales in $P$. Operating profit
is never above $1/4$, so fixed cost two makes staying out strictly optimal.
This verifies entry and off-path pricing on the declared grid. The same
entrant rule and cost draws are used in every comparison.

After each period, an independent public reset occurs with probability
$0<\delta\leq1$ and draws the next state uniformly from $\{C,P\}$.
Without reset, $C$ remains $C$ precisely when both incumbents posted one;
otherwise it moves to $P$. State $P$ remains $P$ without reset. The update
uses public prices and the public reset draw, with no messages. Every joint
price vector satisfies \eqref{eq:minorization} with $\nu=(1/2,1/2)$;
punishment is recurrent rather than absorbing.

The baseline posts one in $C$ and zero in $P$. Each firm's conditional profits
are
\[
 r_C=\frac{1-z}{2}+\frac z3=\frac{3-z}{6},\qquad r_P=0.
\]
The baseline transition has off-diagonal entries $\delta/2$, stationary law
$(1/2,1/2)$ and gain $r_C/2$. Its Poisson equation gives
$h_C-h_P=r_C/\delta$.

\begin{proposition}[Stable elevated prices with entry]\label{prop:entry}
The baseline survives every allowed history-dependent randomized replacement
if and only if
\begin{equation}\label{eq:entrythreshold}
 \delta\leq\frac23-\frac{2z}{9}.
\end{equation}
Its limiting quantity-weighted price exceeds the selected same-market static
benchmark by $1/2$. With positive state-specific demand means $D_C,D_P$ as
specified below, the survival threshold is unchanged and the gap is
$D_C/(D_C+D_P)$.
\end{proposition}
\begin{proof}
Against a baseline rival in $C$, choosing one reproduces the baseline and
has slack zero. Price $3/4$ wins all demand and earns $3/4$; price zero earns
zero. Both cause the nonreset transition to $P$. Substitution in
\eqref{eq:slack} therefore gives
\[
 d(C,3/4)=r_C/\delta-3/4,
 \qquad d(C,0)=r_C/\delta>0.
\]
In $P$, the rival posts zero. Every own price earns zero and has the same
next-state law, so all three slacks are zero. These check every allowed
price, not merely a proposed deviation. Theorem~\ref{thm:replacement} now
gives \eqref{eq:entrythreshold}, including equality and $z=0,1$.
If it fails, choosing $3/4$ in $C$ and zero in $P$ is the theorem's profitable
stationary replacement.

For the benchmark, select both incumbents pricing zero at each state. Against
zero, every grid price earns zero, verifying all stage best-response inequalities.
The entrant uses the identical optimal program; zero incumbent prices give
it no sales. Thus the benchmark transaction price is zero each date. Along
the dynamic baseline all transactions have price one in $C$ and zero in $P$,
and total quantity equals one. Its stationary ratio is $1/2$ from every
initial distribution. The benchmark has both transition rows
$(\delta/2,1-\delta/2)$, unlike the baseline chain, but its transaction price
remains zero. Hence every prespecified material
threshold $0<\varepsilon\leq1/2$ is attained by the limiting gap.

More generally, let demand at state $s$ have a finite strictly positive support
in $(0,1]$, mean $D_s$, and conditional independence of prices, entry cost and
reset given $s$. A uniform support lower bound supplies $q_{\min}>0$.
The entrant's optimization remains valid. Cooperative profit and undercut
profit are both multiplied by $D_C$; the Poisson difference becomes
$r_CD_C/\delta$. Thus the positive factor cancels from the survival threshold.
Expected stationary expenditure is $D_C/2$ and quantity is
$(D_C+D_P)/2$, proving the stated ratio. Demand perturbations preserve
survival and preserve a chosen material threshold exactly when this ratio
stays above it. Changes in $z$ alter the reset threshold by $-2/9$ per unit
change in entry probability. Other entrant behavior changes require recomputing
the finite tables and applying the general test.
\end{proof}

For example, $\delta=1/4$, $z=1/2$ and unit demand yield
$r_C=5/12$, gain $5/24$, $h_C-h_P=5/3$, and undercut slack $11/12$.
There is genuine entry with positive probability and a strict deterrence margin.
The zero slacks in $P$ are structural, since its rival prices zero and all
own prices give the same zero reward and transition.

\begin{proposition}[Persistent exploration can destroy survival]\label{prop:explore}
In this last numerical market, let only incumbent one choose $3/4$ with
probability $e\in(0,1]$ at every $C$ visit and one otherwise; leave the other
incumbent unchanged and both prices zero in $P$. The perturbed profile has
limiting price
\[
 \bar p(e)=\frac{1-e/4}{2(1+3e)}>0,
\]
but replacement by the original deterministic program raises its average
profit by $11e/[24(1+3e)]>0$.
\end{proposition}
\begin{proof}
Put $r=5/12$ and $a=3/4$. The exploring incumbent's reward in $C$ is
$(1-e)r+ea$. Its probability of leaving $C$ is
$\delta/2+(1-\delta)e$ and the $P$-to-$C$ probability is $\delta/2$.
Consequently, with $D=\delta+(1-\delta)e$, its stationary $C$ mass is
$\delta/(2D)$ and its gain is $\delta[(1-e)r+ea]/(2D)$.
Reverting to the original program restores the baseline gain $r/2$.
Subtracting gives exactly
\[
 \frac r2-\frac{\delta[(1-e)r+ea]}{2D}
 =\frac{e(r-\delta a)}{2D}
 =\frac{11e}{24(1+3e)}>0.
\]
All demand in $C$ transacts at one unless exploration occurs, in which case
it transacts at $3/4$. Multiplying conditional mean price $1-e/4$ by the
stationary $C$ mass gives the formula. The benchmark remains zero.
At $e=1/10$ the price gap is $3/8$ and the profitable replacement gain is
$11/312$. As $e\downarrow0$ the gap tends to $1/2$, so any strict material
threshold below $1/2$ survives sufficiently small persistent exploration
even though exact replacement survival fails for every positive rate.
\end{proof}

Finally consider unconditional commitment to price one in this same three-price
market. It earns $r_C$, but constant price $3/4$ earns $3/4>r_C$ against its
rival. Hence elevated unconditional prices fail survival. With a different
declared grid $\{0,1\}$, however, unconditional all-one pricing is a stage
equilibrium: its positive $r_C$ beats a zero-price deviation's zero profit.
All-zero pricing is another valid same-market benchmark selection. This coarse
architecture admits a stable high-price gap without reciprocal punishment.
It is a separate architecture example, not a stability claim against the
three-price replacement set.

\section{Mechanism, stationarity and scope}
Rearranging \eqref{eq:slack} with the baseline Poisson equation gives
\begin{equation}\label{eq:mechanism}
 d_i(s,b)=r_i(s)-r_i(s,b)
       +[P(s,\cdot)-P_i(s,\cdot\mid b)]h_i.
\end{equation}
The first term is the current profit cost of departure. The second is the
continuation effect through observable rival registers, physical demand and
entry. In the example, departure changes the rival's next price through the
public punishment register; that is reciprocal punishment. If rivals commit
to fixed prices while entry or demand responds to current prices, their
continuation effects remain in \eqref{eq:mechanism}; they are not reciprocal
punishment. If the entire continuation kernel is independent of a firm's
price, survival reduces exactly to statewise one-period optimality
$r_i(s,b)\leq r_i(s)$. Thus elevated prices alone establish neither equilibrium
nor communication.

The classification applies to specified finite-register learners and to program
choice within the declared pricing architecture. It determines their limits
and stability but does not assert that an unspecified learning rule discovers
the baseline program. Its stationary environment may include demand or entry
shocks and public regime switches represented by a finite observed state,
provided the enlarged state still satisfies the uniform minorization.
A deterministic cycling phase without phase resets generally violates that
minorization and is not covered merely by finite augmentation.
Arbitrary time trends, unbounded learning counters,
hidden predictive memory, or absorbing punishment without resets fall outside
that architecture. In particular, without full support a negative slack in a
transient state need not yield profitable long-run replacement. The reset
assumption is exactly what makes every one-state switch recurrent in the
necessity proof.

\paragraph{Full-solution scope.}
The target is L13-76 under this explicitly declared finite-grid, finite-public-state
architecture. Lemma~\ref{lem:poisson}, Theorem~\ref{thm:replacement} and
\eqref{eq:gap} cover every increasing integer horizon, fixed initial law and
allowed arbitrary-history replacement. The same-environment selected static
benchmark, entry and denominator requirements are explicit. The perturbation
results cover demand, exploration and entrant behavior, including exact-stability
failure; the example establishes a nonempty stable positive-gap class with
entry. Equation~\eqref{eq:mechanism} and the stated architectural restrictions
address punishment, commitment and stationarity. This settles the original
requests for the specified class, not every possible market or algorithm.

\paragraph{Attribution, sources and verification.}
OpenAI Codex agents generated the mathematical arguments, source checks and
manuscript on 11 October 2026. The platform exposes the GPT-6 family but neither
each agent's exact serving revision nor reasoning setting. JiHo Bae specified
the problem range, preservation, writing and citation requirements and requested
submission; no significant human mathematical contribution or independent
human proof verification is confirmed. All mathematical portions remain
unchecked by an independent human. Separate model reviews, compilation and
repository checks do not constitute independent mathematical verification.
No Lean, simulation or executed controller, optimization or quantifier-elimination
solver is used. The arithmetic examples and finite criteria are proved directly.

The literal repository research Prompt A was applied on 11 October 2026 at
05:39:24 Asia/Seoul, before full-proof construction, to the complete original
statement, contribution policy and actual model facts. Reproducibility notes
retain the input hashes, proof development, source audit and final Prompt B
review record. The exact private generation transcript and unexposed serving
settings are unavailable. The earlier attempt's absorbing no-entry construction
is preserved; it does not prove the present finite criterion, entry threshold
or perturbation classification. Primary Shapiro--Yurukoglu and
Basu--Pollack--Roy texts were consulted only for the roles cited here.
No exhaustive novelty or literature search is claimed.

\section{Completion Estimate}
\noindent\textbf{COMPLETION: 100\%}
For the declared architecture, the proof gives necessary and sufficient conditions
for the full persistence and arbitrary-history replacement target, its
same-environment benchmark, all requested perturbation families and mechanism
and stationarity distinctions. The entry example verifies nonvacuity. This is
a scope claim, not proof confidence or independent verification.

\begin{thebibliography}{9}
\small\raggedright
\bibitem{catalogue}
\emph{Erdosproblems-llm-hunter}, registered economics statement L13-76,
``Algorithmic pricing conduct,'' source identity OP-0398.
\url{https://github.com/mehmetmars7/Erdosproblems-llm-hunter/blob/b1f633121074e16579fa7fd05314e5f504d29bae/attacks/open_problems/economics/statements/L13-76.tex}.
\bibitem{shapiro}
Carl Shapiro and Ali Yurukoglu (2024).
\emph{Trends in Competition in the United States: What Does the Evidence Show?}
Author draft, 24 August 2024, Section 6, p.37, fourth research direction.
\url{https://web.stanford.edu/~ayurukog/trendsincompetition.pdf}.
NBER Working Paper 32762, issued July 2024, revised August 2024:
\url{https://doi.org/10.3386/w32762}.
\bibitem{bpr}
Saugata Basu, Richard Pollack and Marie-Fran\c{c}oise Roy (2006).
\emph{Algorithms in Real Algebraic Geometry}, second edition.
Springer, Algorithms and Computation in Mathematics 10, Sections 13.1 and 14.2.
\url{https://doi.org/10.1007/3-540-33099-2}.
Consulted author-hosted text:
\url{https://www.math.purdue.edu/~sbasu/bpr-posted1.pdf}.
\end{thebibliography}
\end{document}
